% Provenance (never printed): see sources/outline_sources.md (Year 3:
% partition coefficients; batch and fractional melting; crystallisation
% models; reading a spider diagram quantitatively). Brief:
% sources/book5/BRIEFS.md, ch. 12. Curves from figdata/bachelor-3/
% melting_models.py. Exercise and problem values in the Book 5 scratch
% calculator calc/ch12.py. Aldren Island and its lavas are invented.
\chapter{Modelización de elementos en trazas}\label{ch:b3:trace-element-modelling}

Dos basaltos entraron en erupción en la misma isla, separados por pocos kilómetros.
Uno tiene más de una vez y media más lantano que el otro, y casi cuatro veces menos iterbio.
¿Se fundieron a partir de dos mantos distintos, o de un mismo manto fundido en distinta medida,
o a profundidades diferentes? Los elementos mayoritarios apenas permiten distinguirlo: ambas son basaltos.
Los elementos en trazas, que siguen a los minerales sin alterarlos, guardan un registro de cuánto fundido se formó,
de cómo abandonó su fuente y de qué minerales quedaron retenidos.
Este capítulo escribe ese registro en forma de ecuaciones
y lo interpreta a partir de ellas.

\begin{recall}
El coeficiente de reparto mineral--fundido $D$, el coeficiente de reparto global,
los elementos compatibles e incompatibles, las tierras raras, la condrita CI
y el manto primitivo, los diagramas normalizados y la anomalía de europio,
MORB y OIB se encuentran en el volumen de segundo año, en su capítulo sobre los elementos en la
Tierra. La peridotita, la lherzolita, la fusión parcial y la fracción de fundido $F$
están en su capítulo sobre la fusión del manto; la cristalización en equilibrio y la fraccionada,
la ley de fraccionamiento de Rayleigh, la mezcla y la
asimilación--cristalización fraccionada (AFC), en su capítulo sobre la
diferenciación magmática. Allí se demostró la constancia de la razón entre dos elementos muy
incompatibles durante la cristalización.
\end{recall}

\section{El balance de masa de una roca en fusión}\label{sec:b3:trace-element-modelling:balance}

Una roca de masa 1 se funde. A una fracción de fundido $F$ ha generado una masa $F$ de
fundido y conserva una masa $1 - F$ de residuo sólido. Un elemento traza presente en la
fuente a la concentración $C_0$ se reparte entre ambos; su concentración
en el fundido es $C_L$, en el residuo $C_S$. Todo modelo de este capítulo es una
forma de un único balance, $C_0 = FC_L + (1 - F)C_S$ para el conjunto de fundido y residuo,
y una regla que vincula $C_S$ con $C_L$: en el contacto entre un cristal y
el fundido, $C_{\text{mineral}} = D\,C_L$. Los modelos difieren en qué está en contacto
con qué y cuándo.

Los coeficientes de reparto de los minerales del manto difieren mucho. En
una recopilación, el lantano tiene $D = 0.058$ en clinopiroxeno y $0.001$ en
granate, olivino, ortopiroxeno y espinela; el iterbio tiene $D = 0.58$ en
clinopiroxeno, $4.5$ en granate y $0.16$ en los otros tres. El granate incorpora las
tierras raras pesadas y rechaza las ligeras.
% ledger: budahn91:d-ree

\begin{definition}[Fusión no modal]\label{def:b3:trace-element-modelling:nonmodal}
En la \emph{fusión no modal}\index{fusión no modal}, los minerales de la
fuente pasan al fundido en proporciones $p_i$ que difieren de sus fracciones másicas
$x_i$ en la fuente. El fundido ve entonces el coeficiente de reparto global de la
fuente, $D_0 = \sum_i x_iD_i$, al inicio, y el coeficiente
ponderado por la fusión, $P = \sum_i p_iD_i$, a medida que esta avanza.
\end{definition}

La fusión rara vez es modal. En un ejemplo resuelto de lherzolita de espinela con
\qty{58}{\percent} de olivino, \qty{27}{\percent} de ortopiroxeno,
\qty{12}{\percent} de clinopiroxeno y \qty{3}{\percent} de espinela, los minerales
pasan al fundido en las proporciones 20, 25, 45 y \qty{10}{\percent}:
el clinopiroxeno se funde cuatro veces más deprisa que su proporción, y el olivino tres veces
más despacio.
% ledger: white13:ex73

\begin{proposition}[Coeficiente de reparto global del residuo]
\label{prop:b3:trace-element-modelling:residue}
En la \omterm{def:b3:trace-element-modelling:nonmodal}{fusión no modal} con $p_i$ constante, el residuo a la fracción de fundido $F$
tiene el coeficiente de reparto global
\[
\bar D(F) = \frac{D_0 - PF}{1 - F},
\]
válida mientras permanezca cada mineral, es decir, para $F < x_i/p_i$ para cada $i$.
\end{proposition}

\begin{proof}
Por unidad de masa de fuente, la masa del mineral $i$ que queda a $F$ es $x_i - p_iF$,
y el residuo tiene masa $1 - F$. Su coeficiente global es la media ponderada por masa
de los valores, $\sum_i (x_i - p_iF)D_i/(1 - F) = (D_0 - PF)/(1 - F)$. Un mineral se agota
cuando $x_i - p_iF = 0$.
\end{proof}

Con los coeficientes minerales anteriores, la lherzolita de espinela del ejemplo
resuelto y una lherzolita de granate mantenida a \qty{5}{GPa} y
\qty{1400}{\celsius} en el laboratorio, con \qty{54.1}{\percent} de olivino,
\qty{9.8}{\percent} de ortopiroxeno, \qty{19.5}{\percent} de clinopiroxeno y
\qty{16.7}{\percent} de granate en masa, tienen los coeficientes globales siguientes.
% ledger: white13:ex73, rfr25:gkr-5gpa, budahn91:d-ree
\begin{center}
\small
\begin{tabular}{lccc}
\hline
fuente & $D_{\mathrm{La}}$ & $D_{\mathrm{Sm}}$ & $D_{\mathrm{Yb}}$ \\
\hline
lherzolita de espinela & 0.0078 & 0.062 & 0.21 \\
lherzolita de granate & 0.012 & 0.107 & 0.97 \\
\hline
\end{tabular}
\end{center}
El lantano es muy incompatible en ambas. El iterbio es moderadamente incompatible
en la lherzolita de espinela y casi neutro en la lherzolita de granate. La
diferencia entre las dos fuentes reside casi por completo en las tierras raras
pesadas.

\section{Fusión en equilibrio}\label{sec:b3:trace-element-modelling:batch}

\begin{theorem}[Fusión en equilibrio]\label{thm:b3:trace-element-modelling:batch}
En la \emph{fusión en equilibrio}\index{fusión en equilibrio}, el fundido permanece en equilibrio
con la totalidad del residuo hasta salir en un único lote. Con un coeficiente de reparto global
$D$ constante, el fundido y el residuo a la fracción de fundido $F$
tienen
\[
\frac{C_L}{C_0} = \frac{1}{D + F(1 - D)}, \qquad
\frac{C_S}{C_0} = \frac{D}{D + F(1 - D)}.
\]
\end{theorem}

\begin{proof}
En el equilibrio, cada cristal del residuo contiene $D_iC_L$, de modo que el residuo
contiene $C_S = DC_L$ de promedio. El balance $C_0 = FC_L + (1 - F)DC_L$ da
$C_L$, y $C_S$ se obtiene a continuación.
\end{proof}

La fórmula es la ley de cristalización en equilibrio del volumen de segundo año leída
al revés: en ambas, $F$ es la fracción de líquido en un sistema cerrado. Sus dos
límites merecen conocerse. Cuando $F \to 0$, $C_L \to C_0/D$: la primera gota de fundido
puede enriquecer un elemento como máximo en el factor $1/D$. Cuando $D \to 0$, $C_L \to
C_0/F$: un elemento perfectamente incompatible se diluye en el fundido, y su
enriquecimiento depende solo de $F$. Para un elemento compatible, $D > 1$, el fundido
es más pobre que la fuente, y tanto más cuanto menor es $F$.

\begin{proposition}[Fusión en equilibrio no modal]\label{prop:b3:trace-element-modelling:nonmodal}
En la \omterm{thm:b3:trace-element-modelling:batch}{fusión en equilibrio} con las proporciones de fusión $p_i$ de la \cref{def:b3:trace-element-modelling:nonmodal},
\[
\frac{C_L}{C_0} = \frac{1}{D_0 + F(1 - P)}.
\]
\end{proposition}

\begin{proof}
El residuo contiene $\bar D(F)\,C_L$ (la \cref{prop:b3:trace-element-modelling:residue}),
de modo que $C_0 = FC_L + (1 - F)\bar D(F)C_L = C_L(F + D_0 - PF)$.
\end{proof}

Cuando $P$ supera a $D_0$, como ocurre con el lantano en el ejemplo resuelto, donde $D_0 =
0.0074$ y $P = 0.025$, el residuo pierde su mejor huésped, el clinopiroxeno, más deprisa
que el resto. El fundido resulta entonces algo más rico de lo que predice la fusión modal,
en menos de \qty{2}{\percent} a $F = 0.1$ para el lantano (ejercicio~6), pero el
residuo se queda sin clinopiroxeno a $F = 0.12/0.45 \approx 0.27$, y más allá de
ese punto cambian los minerales y los coeficientes.
% ledger: white13:ex73

\section{Fusión fraccionada y dinámica}\label{sec:b3:trace-element-modelling:fractional}

El fundido de una peridotita se acumula a lo largo de los bordes de grano y puede escapar
en cuanto queda conectado. El extremo opuesto a la \omterm{thm:b3:trace-element-modelling:batch}{fusión en equilibrio} elimina cada gota
en cuanto se forma.

\begin{theorem}[Fusión fraccionada]\label{thm:b3:trace-element-modelling:fractional}
En la \emph{fusión fraccionada}\index{fusión fraccionada}, cada incremento de fundido
está en equilibrio con el residuo solo en el momento en que se forma, y después
lo abandona. Con $D$ constante, cuando se ha extraído una masa total $F$ de fundido, el
residuo y el fundido que se forma en ese momento cumplen
\[
\frac{C_S}{C_0} = (1 - F)^{1/D - 1}, \qquad
\frac{C_L}{C_0} = \frac{1}{D}(1 - F)^{1/D - 1}.
\]
\end{theorem}

\begin{proof}
El residuo tiene masa $1 - F$ y concentración $C_S$. Cuando sale una masa
$\mathrm dF$ de fundido a $C_L = C_S/D$, el elemento en el residuo
disminuye en $\mathrm d[(1 - F)C_S] = -(C_S/D)\,\mathrm dF$. Desarrollando,
$(1 - F)\,\mathrm dC_S = C_S(1 - 1/D)\,\mathrm dF$, de modo que
$\mathrm dC_S/C_S = (1/D - 1)\,\mathrm d[-\ln(1 - F)]$. Integrando desde
$C_S = C_0$ en $F = 0$ resulta $\ln(C_S/C_0) = (1/D - 1)\ln(1 - F)$, y
$C_L = C_S/D$.
\end{proof}

Para un elemento incompatible, $1/D - 1$ es grande y positivo: el residuo queda
empobrecido en pocas veces $D$ de fusión, y el fundido instantáneo cae
desde $C_0/D$ hasta casi nada. Un residuo de \omterm{thm:b3:trace-element-modelling:fractional}{fusión fraccionada} está mucho más
empobrecido que uno de \omterm{thm:b3:trace-element-modelling:batch}{fusión en equilibrio}. Los fundidos, sin embargo, no se extraen
uno a uno: se acumulan en canales y se reúnen en charcos, y el \omterm{def:b3:trace-element-modelling:fractional}{fundido fraccionado agregado},
definido más abajo, se dibuja con trazos en la figura.

\begin{omfigure}
\begin{tikzpicture}
\begin{groupplot}[group style={group size=2 by 1, horizontal sep=1.6cm},
  omrecord, width=0.47\linewidth, height=0.44\linewidth, title style={font=\small},
  xlabel={fracción de fundido $F$}, xmin=0, xmax=0.3, ymode=log, ymin=0.01, ymax=150,
  xtick={0,0.1,0.2,0.3}, xticklabels={0,0.1,0.2,0.3},
  ytick={0.01,0.1,1,10,100}, yticklabels={0.01,0.1,1,10,100},
  legend style={font=\scriptsize, legend cell align=left}]
\nextgroupplot[title={(a) en equilibrio (continua), agregada (de trazos)},
  ylabel={$C_L/C_0$}, legend style={at={(0.97,0.97)}, anchor=north east}]
% computed in figdata/bachelor-3/melting_models.py
\addplot[omDef, thick] table[x=f, y=b001] {figdata/out/bachelor-3-melting_models-curves.dat};
\addlegendentry{$D = 0.01$}
\addplot[red!70!black, thick] table[x=f, y=b01] {figdata/out/bachelor-3-melting_models-curves.dat};
\addlegendentry{$D = 0.1$}
\addplot[black, thick] table[x=f, y=b1] {figdata/out/bachelor-3-melting_models-curves.dat};
\addlegendentry{$D = 1$}
\addplot[omLithLime!70!black, thick] table[x=f, y=b10] {figdata/out/bachelor-3-melting_models-curves.dat};
\addlegendentry{$D = 10$}
\addplot[omDef, thick, densely dashed] table[x=f, y=a001] {figdata/out/bachelor-3-melting_models-curves.dat};
\addplot[red!70!black, thick, densely dashed] table[x=f, y=a01] {figdata/out/bachelor-3-melting_models-curves.dat};
\addplot[omLithLime!70!black, thick, densely dashed] table[x=f, y=a10] {figdata/out/bachelor-3-melting_models-curves.dat};
\nextgroupplot[title={(b) fraccionada, instantánea}]
\addplot[omDef, thick] table[x=f, y=i001] {figdata/out/bachelor-3-melting_models-curves.dat};
\addplot[red!70!black, thick] table[x=f, y=i01] {figdata/out/bachelor-3-melting_models-curves.dat};
\addplot[black, thick] table[x=f, y=i1] {figdata/out/bachelor-3-melting_models-curves.dat};
\addplot[omLithLime!70!black, thick] table[x=f, y=i10] {figdata/out/bachelor-3-melting_models-curves.dat};
\end{groupplot}
\end{tikzpicture}

\par\smallskip{\small Enriquecimiento de un elemento en trazas en el fundido frente a la
fracción de fundido, para cuatro coeficientes de reparto global. (a) Fundidos en equilibrio (continua)
y \omterm{def:b3:trace-element-modelling:fractional}{fundidos fraccionados agregados} (de trazos; para $D = 1$ ambos permanecen en 1). (b) Los
fundidos instantáneos de la \omterm{thm:b3:trace-element-modelling:fractional}{fusión fraccionada}: un elemento con $D = 0.01$ se
agota del residuo tras unos pocos por ciento de fusión.}\label{fig:b3:trace-element-modelling:curves}
\end{omfigure}

\begin{definition}[Fundido fraccionado agregado]\label{def:b3:trace-element-modelling:fractional}
El \emph{fundido fraccionado agregado}\index{fundido fraccionado agregado} es la
mezcla de todos los fundidos instantáneos de una \omterm{thm:b3:trace-element-modelling:fractional}{fusión fraccionada} desde $0$ hasta
$F$, acumulados sin reacción posterior con el residuo.
\end{definition}

\begin{proposition}[Composición del fundido agregado]
\label{prop:b3:trace-element-modelling:aggregate}
Con un $D$ constante, el \omterm{def:b3:trace-element-modelling:fractional}{fundido fraccionado agregado} cumple
\[
\frac{\bar C_L}{C_0} = \frac{1 - (1 - F)^{1/D}}{F}.
\]
\end{proposition}

\begin{proof}
El elemento extraído es la fuente menos el residuo,
$C_0 - (1 - F)C_S = C_0[1 - (1 - F)^{1/D}]$ según el \cref{thm:b3:trace-element-modelling:fractional},
repartido sobre una masa $F$ de fundido.
\end{proof}

Cuando $(1 - F)^{1/D}$ es pequeño, el fundido agregado contiene todo el
elemento y $\bar C_L \approx C_0/F$, sea cual sea $D$: para elementos muy
incompatibles, los \omterm{def:b3:trace-element-modelling:fractional}{fundidos fraccionados agregados} y los fundidos en equilibrio
con $D$ pequeño coinciden, y ambos miden $F$.

La fusión real se sitúa entre los dos extremos. El residuo retiene algo de
fundido en sus poros, y el fundido que sale es el exceso por encima de esa porosidad.

\begin{definition}[Fusión dinámica]\label{def:b3:trace-element-modelling:dynamic}
En la \emph{fusión dinámica}\index{fusión dinámica} (también llamada fusión
continua), el residuo conserva una fracción másica constante $\phi$ de fundido en
sus poros, en equilibrio con sus cristales, y todo el fundido producido más allá de
esa cantidad se extrae.
\end{definition}

\begin{proposition}[Fusión dinámica]\label{prop:b3:trace-element-modelling:dynamic}
La \omterm{def:b3:trace-element-modelling:dynamic}{fusión dinámica} con porosidad $\phi$ es una \omterm{thm:b3:trace-element-modelling:fractional}{fusión fraccionada} con el
coeficiente de reparto efectivo $D' = (1 - \phi)D + \phi$, donde $F$ es la
fracción másica de fundido extraída:
\[
\frac{C_L}{C_0} = \frac{1}{D'}(1 - F)^{1/D' - 1}.
\]
\end{proposition}

\begin{proof}
El residuo está formado por cristales con la fracción $1 - \phi$ y por fundido
de poros con la fracción $\phi$, ambos en equilibrio con el fundido a $C_L$, de
modo que $C_S = [(1 - \phi)D + \phi]C_L = D'C_L$. La demostración
del \cref{thm:b3:trace-element-modelling:fractional} se aplica entonces con $D'$.
\end{proof}

Para $\phi = 0$ el modelo es la \omterm{thm:b3:trace-element-modelling:fractional}{fusión fraccionada}; para $\phi \to 1$ se aproxima
a la \omterm{thm:b3:trace-element-modelling:batch}{fusión en equilibrio}. Una porosidad del uno por ciento limita el
enriquecimiento del primer fundido a $1/\phi = 100$, por pequeño que sea $D$: el
fundido de poros amortigua el residuo.

\begin{omfigure}
\begin{tikzpicture}[font=\small, >=Stealth]
\tikzset{rock/.style={draw, fill=omLithLime!25, minimum width=2.3cm, minimum height=2.0cm},
  melt/.style={draw, fill=red!25}}
% (a) batch
\node[rock] (a) at (0,0) {};
\foreach \x/\y in {-0.7/0.5, 0.3/0.6, -0.2/-0.1, 0.7/-0.4, -0.6/-0.6, 0.4/0.0}
  \fill[red!55] (\x,\y) circle (0.09);
\node[align=center, font=\scriptsize] at (0,-1.45) {residuo $1-F$ con $DC_L$\\fundido $F$ con $C_L$};
\draw[->, thick, red!70!black] (0,1.05) -- (0,1.7) node[above, font=\scriptsize] {todo de una vez};
\node[font=\small] at (0,-2.25) {(a) en equilibrio};
% (b) fractional
\node[rock] (b) at (4,0) {};
\draw[->, thick, red!70!black] (4,1.05) -- (4,1.7);
\foreach \y/\s in {1.85/0.06, 2.05/0.08, 2.27/0.1}
  \fill[red!55] (4,\y) circle (\s);
\node[align=center, font=\scriptsize] at (4,-1.45) {cada gota sale\\al formarse};
\node[font=\small] at (4,-2.25) {(b) fraccionada};
% (c) dynamic
\node[rock] (c) at (8,0) {};
\foreach \x/\y in {7.4/0.5, 8.4/0.3, 7.8/-0.5, 8.6/-0.6}
  \fill[red!55] (\x,\y) circle (0.06);
\draw[->, thick, red!70!black] (8,1.05) -- (8,1.7);
\foreach \y/\s in {1.85/0.06, 2.05/0.08, 2.27/0.1}
  \fill[red!55] (8,\y) circle (\s);
\node[align=center, font=\scriptsize] at (8,-1.45) {los poros retienen $\phi$;\\sale el exceso};
\node[font=\small] at (8,-2.25) {(c) dinámica};
\end{tikzpicture}

\par\smallskip{\small Tres formas de que el fundido salga de una fuente en fusión. (a)
En equilibrio: el fundido permanece con todo el residuo y sale después. (b) Fraccionada:
cada incremento sale de inmediato; los incrementos acumulados forman el fundido
agregado. (c) Dinámica: el residuo conserva una pequeña fracción másica $\phi$ de
fundido en sus poros.}\label{fig:b3:trace-element-modelling:modes}
\end{omfigure}

\section{Fusión y cristalización}\label{sec:b3:trace-element-modelling:crystallisation}

Un fundido que sale de su fuente rara vez llega a la superficie sin cambios. Va
cristalizando olivino, luego piroxeno y plagioclasa durante el ascenso, puede
asimilar roca de caja y puede mezclarse con otros lotes. Las leyes del volumen de
segundo año se aplican por turnos: una lava es un fundido de fracción $F$ de su
fuente, y después un líquido que ha conservado una fracción $f$ de su masa mediante
cristalización fraccionada, de modo que, para un elemento con $D$ en el residuo de
fusión y $D_c$ en los minerales que cristalizan,
\[
C_{\text{lava}} = C_0\,\frac{1}{D + F(1 - D)}\,f^{\,D_c - 1}.
\]
El segundo factor es el mismo para todos los elementos que los minerales en
cristalización rechazan, $D_c \approx 0$: los escala a todos por $1/f$. Los cocientes entre
tales elementos quedan, por tanto, fijados en la fuente, y una lava que ha perdido olivino
conserva el cociente La/Yb de su fundido primario mientras su lantano aumenta. Invertir $F$
a partir de una concentración requiere también $f$; invertir la fuente a partir de un cociente
no lo requiere. La plagioclasa es la excepción entre las tierras raras: incorpora europio,
y su eliminación deja la anomalía de europio del volumen de segundo año.

\section{Interpretación de los patrones de tierras raras}\label{sec:b3:trace-element-modelling:patterns}

Un fundido de una lherzolita de espinela y un fundido de una lherzolita de granate
comparten sus tierras raras ligeras, que ningún residuo retiene, y difieren en las
pesadas, que retiene el granate. En un diagrama normalizado a condrita, los fundidos de
granate son empinados y de fondo plano en el extremo pesado; los de espinela son más suaves. La
pendiente aumenta a medida que $F$ disminuye.

\begin{omfigure}
\begin{tikzpicture}
\begin{axis}[omrecord, width=0.8\linewidth, height=0.5\linewidth, ymode=log,
  xmin=56.5, xmax=71.5, ymin=1, ymax=1000, xtick={57,58,60,62,63,64,65,70,71},
  xticklabels={La,Ce,Nd,Sm,Eu,Gd,Tb,Yb,Lu}, ytick={1,10,100,1000}, yticklabels={1,10,100,1000},
  ylabel={muestra / condrita CI},
  legend style={font=\scriptsize, at={(0.98,0.98)}, anchor=north east,
  legend cell align=left, legend columns=2}]
% computed in figdata/bachelor-3/melting_models.py
\addplot[omDef, thick, mark=*, mark size=1.3pt] table[x=z, y=sp01] {figdata/out/bachelor-3-melting_models-ree.dat};
\addlegendentry{espinela, $F = 1\,\%$}
\addplot[red!70!black, thick, mark=square*, mark size=1.3pt] table[x=z, y=gt01] {figdata/out/bachelor-3-melting_models-ree.dat};
\addlegendentry{granate, $F = 1\,\%$}
\addplot[omDef, thick, densely dashed, mark=*, mark size=1.3pt] table[x=z, y=sp05] {figdata/out/bachelor-3-melting_models-ree.dat};
\addlegendentry{espinela, $5\,\%$}
\addplot[red!70!black, thick, densely dashed, mark=square*, mark size=1.3pt] table[x=z, y=gt05] {figdata/out/bachelor-3-melting_models-ree.dat};
\addlegendentry{granate, $5\,\%$}
\addplot[omDef, thick, densely dotted, mark=*, mark size=1.3pt] table[x=z, y=sp20] {figdata/out/bachelor-3-melting_models-ree.dat};
\addlegendentry{espinela, $20\,\%$}
\addplot[red!70!black, thick, densely dotted, mark=square*, mark size=1.3pt] table[x=z, y=gt20] {figdata/out/bachelor-3-melting_models-ree.dat};
\addlegendentry{granate, $20\,\%$}
\addplot[black, thin, mark=none] table[x=z, y=pm] {figdata/out/bachelor-3-melting_models-ree.dat};
\addlegendentry{fuente}
\end{axis}
\end{tikzpicture}

\par\smallskip{\small Patrones de tierras raras normalizados a condrita de fundidos de \omterm{thm:b3:trace-element-modelling:batch}{fusión en equilibrio}
de una fuente de manto primitivo, a partir de una lherzolita de espinela (círculos azules) y una
lherzolita de granate (cuadrados rojos), a tres fracciones de fundido; la fuente es la
línea fina. Solo se representan las nueve tierras raras con coeficientes de reparto
tabulados.}\label{fig:b3:trace-element-modelling:ree}
\end{omfigure}
% ledger: budahn91:d-ree, ms95:pm-trace, lodders:ci-table, white13:ex73, rfr25:gkr-5gpa

El granate cede el paso a la espinela como mineral aluminoso del manto, entre unos 60
y \qty{90}{km} de profundidad. Un fundido cuyas tierras raras pesadas son planas y bajas, con una
pendiente pronunciada de ligeras a pesadas, ha dejado, por tanto, granate en su residuo y
ha empezado a fundirse al menos a esa profundidad.
% ledger: white13:sp-gt

Un patrón mezcla dos incógnitas, $F$ y la fuente. Representar dos elementos
uno frente a otro permite separarlas.

\begin{definition}[Diagrama razón--razón]\label{def:b3:trace-element-modelling:ratios}
Un \emph{diagrama razón--razón}\index{diagrama razón--razón} representa, para una serie
de lavas, la razón entre un elemento muy incompatible $H$ y otro menos
incompatible $M$ frente a la concentración de $H$ (o frente a otra
razón). La fusión parcial, la cristalización y la mezcla dibujan líneas de distintas
formas en él.
\end{definition}

\begin{proposition}[La recta de fusión en equilibrio]\label{prop:b3:trace-element-modelling:line}
Para fundidos de \omterm{thm:b3:trace-element-modelling:batch}{fusión en equilibrio} de una fuente con $D_H = 0$ y un $D_M$ constante, los puntos
$(C_H, C_H/C_M)$ se sitúan sobre la recta
\[
\frac{C_H}{C_M} = \frac{D_M}{C_0^M}\,C_H + \frac{C_0^H(1 - D_M)}{C_0^M}.
\]
\end{proposition}

\begin{proof}
\sloppy Con $D_H = 0$, $C_H = C_0^H/F$, luego $1/F = C_H/C_0^H$. Entonces
$C_H/C_M = (C_0^H/F)\,(D_M + F(1 - D_M))/C_0^M = C_0^HD_M/(FC_0^M) +
C_0^H(1 - D_M)/C_0^M$, y el primer término es $(D_M/C_0^M)\,C_H$.
\end{proof}

La cristalización fraccionada de minerales que rechazan ambos elementos desplaza una lava
horizontalmente en este diagrama, ya que la razón no cambia; la fusión la desplaza
a lo largo de una recta inclinada cuya pendiente crece con $D_M$; la mezcla de dos fundidos dibuja una
hipérbola. Una fuente con granate tiene un $D_{\mathrm{Yb}}$ grande y, por tanto, una
recta de fusión empinada en La/Yb frente a La.

\begin{method}[Invertir dos lavas]\label{met:b3:trace-element-modelling:invert}
(1) Se comprueba que las lavas son cogenéticas y que $H$ es mucho más incompatible
que $M$. (2) Se corrige cada lava por cristalización, o se usan solo razones.
(3) Con dos o más lavas, se ajusta la recta $C_H/C_M = aC_H + b$. (4) Entonces
$a = D_M/C_0^M$ y $b = C_0^H(1 - D_M)/C_0^M$: suponiendo conocido $C_0^M$, $D_M =
aC_0^M$ y $C_0^H = bC_0^M/(1 - D_M)$. (5) La fracción de fundido de cada lava es
$F = C_0^H/C_H$.
\end{method}

\begin{omfigure}
\begin{tikzpicture}
\begin{axis}[omrecord, width=0.78\linewidth, height=0.52\linewidth,
  xmode=log, ymode=log, xmin=2, xmax=60, ymin=2, ymax=100,
  xlabel={La (\unit{ppm})}, ylabel={La/Yb},
  xtick={2,5,10,20,50}, xticklabels={2,5,10,20,50},
  ytick={2,5,10,20,50,100}, yticklabels={2,5,10,20,50,100},
  legend style={font=\scriptsize, at={(0.03,0.97)}, anchor=north west,
  legend cell align=left}]
% computed in figdata/bachelor-3/melting_models.py
\addplot[omDef, thick] table[x=la_sp, y=lay_sp] {figdata/out/bachelor-3-melting_models-traj.dat};
\addlegendentry{lherzolita de espinela}
\addplot[red!70!black, thick] table[x=la_gt, y=lay_gt] {figdata/out/bachelor-3-melting_models-traj.dat};
\addlegendentry{lherzolita de granate}
\addplot[omDef, only marks, mark=|, mark size=3pt, forget plot] table[x=la_sp, y=lay_sp] {figdata/out/bachelor-3-melting_models-ticks.dat};
\addplot[red!70!black, only marks, mark=|, mark size=3pt, forget plot] table[x=la_gt, y=lay_gt] {figdata/out/bachelor-3-melting_models-ticks.dat};
\addplot[black, only marks, mark=diamond*, mark size=2.8pt] table[x=la, y=layb] {figdata/out/bachelor-3-melting_models-lavas.dat};
\addlegendentry{lavas A y B}
\node[font=\scriptsize, anchor=north] at (axis cs:9.6,5.2) {A};
\node[font=\scriptsize, anchor=south east] at (axis cs:15.2,36) {B};
\node[font=\scriptsize, omDef, anchor=north west] at (axis cs:36.5,17.5) {1\,\%};
\node[font=\scriptsize, omDef, anchor=north west] at (axis cs:11.3,6.2) {5\,\%};
\node[font=\scriptsize, omDef, anchor=north west] at (axis cs:3.14,2.55) {20\,\%};
\node[font=\scriptsize, red!70!black, anchor=south east] at (axis cs:29.5,66) {1\,\%};
\node[font=\scriptsize, red!70!black, anchor=south east] at (axis cs:10.5,24) {5\,\%};
\node[font=\scriptsize, red!70!black, anchor=south east] at (axis cs:3.09,7.0) {20\,\%};
\end{axis}
\end{tikzpicture}

\par\smallskip{\small Fundidos de \omterm{thm:b3:trace-element-modelling:batch}{fusión en equilibrio} de una fuente de manto primitivo en un diagrama de La/Yb
frente a La, para una lherzolita de espinela y una de granate, con marcas en $F = 1$,
2, 5, 10 y \qty{20}{\percent}, y las dos lavas inventadas del problema de fin
de semana. Cada lava se sitúa sobre una sola curva.}\label{fig:b3:trace-element-modelling:ratios}
\end{omfigure}

En escala logarítmica la recta de la proposición se curva, y con
$D_{\mathrm{La}} \ne 0$ solo es casi recta en ejes lineales; las curvas
de la figura son exactas. Una fuente de espinela no puede producir
un fundido con razón La/Yb superior a $(C_0^{\mathrm{La}}/D_{\mathrm{La}})/(C_0^{\mathrm{Yb}}/D_{\mathrm{Yb}})$,
unas 40 para el manto primitivo; una fuente de granate puede alcanzar unas 120.

\section{Ejercicios}

\begin{exercise}[$\star$]\label{exo:b3:trace-element-modelling:1}
Con los coeficientes minerales del samario, 0.40 en clinopiroxeno, 0.11 en
granate y 0.016 en olivino, ortopiroxeno y espinela, calcúlese el
coeficiente de reparto global del samario en la lherzolita de espinela (\qty{12}{\percent}
de clinopiroxeno) y en la lherzolita de granate (\qty{19.5}{\percent} de clinopiroxeno,
\qty{16.7}{\percent} de granate) del capítulo.
\end{exercise}
% ledger: budahn91:d-ree

\begin{exercise}[$\star$]\label{exo:b3:trace-element-modelling:2}
Un elemento con $D = 0.05$ está presente a \qty{1}{ppm} en una fuente. Calcúlese
su concentración en los fundidos de \omterm{thm:b3:trace-element-modelling:batch}{fusión en equilibrio} para $F = 0.02$ y $0.10$.
\end{exercise}

\begin{exercise}[$\star$]\label{exo:b3:trace-element-modelling:3}
¿Cuál es el mayor enriquecimiento $C_L/C_0$ que la \omterm{thm:b3:trace-element-modelling:batch}{fusión en equilibrio} puede dar a un
elemento con $D = 0.01$? ¿Y a un elemento con $D = 0$ para
$F = \qty{5}{\percent}$?
\end{exercise}

\begin{exercise}[$\star$]\label{exo:b3:trace-element-modelling:4}
Una lava contiene \qty{30}{ppm} de La y \qty{2.5}{ppm} de Yb. Con los valores de la
condrita CI, La 0.232 e Yb \qty{0.163}{ppm}, calcúlese su razón La/Yb y su
$(\mathrm{La}/\mathrm{Yb})_N$ normalizada a condrita.
\end{exercise}
% ledger: lodders:ci-table

\begin{exercise}[$\star\star$]\label{exo:b3:trace-element-modelling:5}
Para un elemento con $D = 0.01$, compárense los residuos tras un \qty{5}{\percent}
de \omterm{thm:b3:trace-element-modelling:batch}{fusión en equilibrio} y de \omterm{thm:b3:trace-element-modelling:fractional}{fusión fraccionada}. ¿Qué contiene el fundido fraccionado
instantáneo en ese momento?
\end{exercise}

\begin{exercise}[$\star\star$]\label{exo:b3:trace-element-modelling:6}
En el ejemplo resuelto del capítulo (modas 58, 27, 12, \qty{3}{\percent};
proporciones de fusión 20, 25, 45, \qty{10}{\percent} para olivino,
ortopiroxeno, clinopiroxeno y espinela), el lantano tiene los coeficientes
0.0003, 0.002, 0.053 y 0.01. Calcúlense $D_0$ y $P$, y después el enriquecimiento del
lantano en el fundido de \omterm{thm:b3:trace-element-modelling:batch}{fusión en equilibrio} para $F = 0.02$ y $0.10$, no modal y modal.
\end{exercise}
% ledger: white13:ex73

\begin{exercise}[$\star\star$]\label{exo:b3:trace-element-modelling:7}
Para $F = 0.10$, compárense el fundido de \omterm{thm:b3:trace-element-modelling:batch}{fusión en equilibrio} y el \omterm{def:b3:trace-element-modelling:fractional}{fundido fraccionado
agregado} para $D = 0.1$ y para $D = 0.01$. ¿Cuál de los dos difiere más, y por qué?
\end{exercise}

\begin{exercise}[$\star\star$]\label{exo:b3:trace-element-modelling:8}
Un elemento con $D = 0.001$ se funde dinámicamente con $\phi = 0.01$. Calcúlense
$D'$, el enriquecimiento del primer fundido extraído y el del fundido extraído
a $F = 0.01$; compárese con la \omterm{thm:b3:trace-element-modelling:fractional}{fusión fraccionada}.
\end{exercise}

\begin{exercise}[$\star\star$]\label{exo:b3:trace-element-modelling:9}
En la lherzolita de espinela del capítulo, $D_{\mathrm{La}} = 0.0078$ y
$D_{\mathrm{Sm}} = 0.062$. ¿En qué factor la \omterm{thm:b3:trace-element-modelling:batch}{fusión en equilibrio} eleva la razón La/Sm
por encima del valor de la fuente para $F = 0.01$ y para $F = 0.10$? ¿Qué indica eso
sobre el uso de razones para estimar $F$?
\end{exercise}

\begin{exercise}[$\star\star\star$]\label{exo:b3:trace-element-modelling:10}
Demuéstrese que el \omterm{def:b3:trace-element-modelling:fractional}{fundido fraccionado agregado} es la media de los fundidos
instantáneos desde 0 hasta $F$, y que tiende a $C_0/F$ cuando $(1 - F)^{1/D} \to 0$.
Para $D = 0.01$, ¿a partir de qué $F$ está dentro del \qty{1}{\percent} de $C_0/F$?
\end{exercise}

\begin{exercise}[$\star\star\star$]\label{exo:b3:trace-element-modelling:11}
Dos lavas cogenéticas P y Q tienen La de 20 y \qty{10}{ppm}, y La/Yb de 10 y 6.
Suponiendo el lantano perfectamente incompatible y una fuente con
\qty{0.44}{ppm} de Yb, aplíquese el método del capítulo: calcúlense $D_{\mathrm{Yb}}$,
el La de la fuente y la fracción de fusión de cada lava.
\end{exercise}

\begin{exercise}[$\star\star\star$]\label{exo:b3:trace-element-modelling:12}
En la \omterm{def:b3:trace-element-modelling:nonmodal}{fusión no modal} del ejercicio~6, ¿a qué fracción de fusión
se agota cada mineral? Calcúlese el coeficiente global del residuo para el lantano en
$F = 0.2$, y explíquese por qué la fórmula del capítulo no puede usarse más allá del
primer agotamiento.
\end{exercise}

\section{Problema: dos lavas de la isla de Aldren}

\begin{problem}[{Problema de fin de semana --- la fuente, la profundidad y el grado de fusión
de dos basaltos inventados, a partir de su lantano e iterbio}]\label{pb:b3:trace-element-modelling:1}
En la isla de Aldren, una isla volcánica inventada, la lava A contiene \qty{9.6}{ppm} de La
y \qty{1.7}{ppm} de Yb, la lava B \qty{15.6}{ppm} de La y \qty{0.45}{ppm} de Yb.
Ambas se toman como fundidos primarios de una fuente de manto primitivo con
\qty{0.648}{ppm} de La y \qty{0.441}{ppm} de Yb (condrita CI: La 0.232, Yb
\qty{0.163}{ppm}). Coeficientes minerales: La 0.058 en clinopiroxeno y 0.001
en granate y en olivino, ortopiroxeno y espinela; Yb 0.58 en
clinopiroxeno, 4.5 en granate y 0.16 en olivino, ortopiroxeno y espinela.
La lherzolita de espinela tiene \qty{58}{\percent} de olivino, \qty{27}{\percent} de ortopiroxeno, \qty{12}{\percent}
de clinopiroxeno, \qty{3}{\percent} de espinela; la lherzolita de granate \qty{54.1}{\percent} de olivino, \qty{9.8}{\percent}
de ortopiroxeno, \qty{19.5}{\percent} de clinopiroxeno, \qty{16.7}{\percent} de granate. La fusión es modal.
% ledger: ms95:pm-trace, lodders:ci-table, budahn91:d-ree, white13:ex73, rfr25:gkr-5gpa

\medskip\noindent\textbf{Parte I --- Las dos fuentes.}
\begin{enumerate}
\item Calcúlense $D_{\mathrm{La}}$ y $D_{\mathrm{Yb}}$ de la lherzolita de espinela.
\item Calcúlense para la lherzolita de granate. Redondéense los cuatro a dos cifras significativas
cifras significativas para el resto del problema.
\item ¿Qué coeficiente difiere más entre las fuentes, y por qué? ¿Es
el iterbio compatible en la lherzolita de granate?
\item Calcúlese el $(\mathrm{La}/\mathrm{Yb})_N$ normalizado a condrita de la
fuente. ¿Qué aspecto tiene su patrón?
\end{enumerate}

\medskip\noindent\textbf{Parte II --- Fusión en equilibrio del lava A.}
\begin{enumerate}[resume]
\item Calcúlese el La en fundidos en equilibrio de la lherzolita de espinela para $F = 1$, 5 y
\qty{10}{\percent}.
\item Búsquese la fracción de fundido que da al lava A su lantano.
\item Prediganse su iterbio y su La/Yb, y compárense con los del lava A.
\item Calcúlense el La y el Yb en el residuo, y la fracción del
lantano de la fuente que queda en él.
\item ¿Qué fracción de fundido se obtendría si se tomara el lantano como perfectamente
incompatible?
\end{enumerate}

\medskip\noindent\textbf{Parte III --- Fusión fraccionada del lava A.}
\begin{enumerate}[resume]
\item Calcúlese el La en el fundido fraccionado instantáneo para $F = 0.1$, 1 y
\qty{6}{\percent}.
\item Calcúlese el La en el \omterm{def:b3:trace-element-modelling:fractional}{fundido fraccionado agregado} para $F = \qty{6.0}{\percent}$.
\item Búsquese la fracción de fundido a la que el fundido agregado tiene el lantano del
lava A. ¿Por qué es tan próxima a la respuesta de la pregunta~9?
\item Compárese el lantano del residuo con el de la pregunta~8.
\item ¿Qué estimación de $F$ se conserva?
\end{enumerate}

\medskip\noindent\textbf{Parte IV --- Lava B y granate.}
\begin{enumerate}[resume]
\item Calcúlense La/Yb y $(\mathrm{La}/\mathrm{Yb})_N$ de ambas lavas.
\item ¿Cuál es el mayor La/Yb que puede
alcanzar un fundido en equilibrio de la lherzolita de espinela?
\item ¿A qué fracción de fundido tendría un fundido de lherzolita de espinela el La/Yb del
lava B, y cuánto lantano contendría? Concluyase.
\item Búsquese la fracción de fundido del lava B a partir de su lantano, para la
lherzolita de granate, y prediganse su iterbio.
\item ¿Qué dicen las dos fuentes sobre las profundidades a las que empezaron a
fundirse A y B? ¿Podría el lava A proceder de la lherzolita de granate?
\item Calcúlense la pendiente y la ordenada en el origen de la recta de \omterm{thm:b3:trace-element-modelling:batch}{fusión en equilibrio} La/Yb $= a\,
\mathrm{La} + b$ de la fuente de granate, con $D_{\mathrm{La}}$ tomado como cero,
y compárese el lava B con ella.
\item Si el lava A hubiera perdido el \qty{10}{\percent} de su masa como olivino antes de
erupcionar, ¿cómo cambiarían su La/Yb y su fracción de fundido?
\item Respóndase a la pregunta de la apertura del capítulo para estas dos lavas.
\item Indíquese el grado de fusión del lava A.
\end{enumerate}
\end{problem}
